\documentclass[11pt,a4paper]{article}
\usepackage[margin=1in]{geometry}
\usepackage{amsmath,amssymb}
\usepackage{booktabs}
\usepackage{hyperref}

\newcommand{\dist}{\operatorname{d}}
\newcommand{\pos}{\operatorname{pos}}

\title{Exact Domination of Lecture Timings}
\author{Devesh Kumar (bmat2319), Arjina Jana (bmat2312)}
\date{}

\begin{document}
\maketitle

\noindent\textbf{Code:} \url{https://github.com/deveshk-boop/daa-project-q17}
(build and run instructions are in the README).

\section{The problem}

We have $n$ lectures, one per day, each with a start and end time (possibly
crossing midnight). We must pick a subset $S$ of lectures so that every lecture
overlaps \emph{exactly one} lecture of $S$; a chosen lecture counts as
overlapping itself. If no such subset exists, we say so.

\paragraph{Setup.} Times are in minutes on a circle of length $M=1440$. Lecture
$i$ is the arc $A_i$ of $\ell_i=(e_i-s_i)\bmod M$ minutes starting at $s_i$. I
made two assumptions since the statement is silent on them: lectures are
half-open, so back-to-back lectures do not overlap, and a lecture with
$s_i=e_i$ lasts one minute. Write $\dist(x,y)=(y-x)\bmod M$ for the clockwise
distance from $x$ to $y$; a time $p$ is in $A_i$ iff $\dist(s_i,p)<\ell_i$.

\section{Key observations}

\textbf{Overlap test.} Two arcs overlap iff one of them starts inside the other.
(If $p$ is in both, the one whose start is closer behind $p$ starts inside the
other.) So the overlap table takes $O(1)$ per pair.

\textbf{Chosen lectures are disjoint.} A chosen lecture already overlaps itself,
so it may not overlap any other chosen lecture.

\textbf{Single lecture.} $S=\{a\}$ is valid iff $a$ overlaps every lecture.

\textbf{Two or more lectures.} Then $S$ is a set of disjoint arcs, so they sit
around the circle in a fixed order $x_0,x_1,\dots,x_{k-1}$. For two neighbours
$a$ then $b$, the \emph{gap} is the time between the end of $a$ and the start of
$b$. Say $b$ \emph{can follow} $a$, written $F(a,b)$, if:
\begin{enumerate}
  \item $a$ and $b$ do not overlap;
  \item no lecture overlaps both $a$ and $b$ (it would be covered twice);
  \item no lecture that overlaps neither $a$ nor $b$ starts inside the gap
  (it would not be covered at all).
\end{enumerate}

\paragraph{Claim.} A set $S$ of $k\ge2$ disjoint arcs is valid iff $F(x_i,x_{i+1})$
holds for every neighbouring pair (cyclically, including $x_{k-1}\to x_0$).

\emph{Proof.} ($\Rightarrow$) If $S$ is valid, condition 1 holds because the
arcs are disjoint. Condition 2 holds, otherwise some lecture is covered twice.
For condition 3, suppose $c$ overlaps neither $a$ nor $b$ but starts in their
gap. It cannot reach the start of $b$, so it lies entirely inside the gap. No
chosen arc touches the gap, because $a$ and $b$ are neighbours, so $c$ overlaps
nothing in $S$, a contradiction.

($\Leftarrow$) Take any lecture $c$. If $c$ overlapped two chosen arcs, then, since
$c$ is one connected stretch of the circle, walking from one of them to the other
it must touch a pair of neighbours, and condition 2 forbids that. If $c$
overlapped no chosen arc, its start would fall in some gap, and it would overlap
neither neighbour of that gap, which condition 3 forbids. So $c$
overlaps exactly one chosen arc. $\square$

\section{Algorithm}

\begin{enumerate}
  \item Build the overlap table.
  \item If some lecture overlaps all lectures, output it.
  \item Otherwise build the table $F(a,b)$ for all pairs, by checking the three
  conditions against every lecture $c$.
  \item Lecture~1 must overlap some chosen lecture, so try each lecture $a_0$
  that overlaps lecture~1 as the starting point. For each one:
  \begin{itemize}
    \item Sort the other lectures by $\pos(x)=\dist(s_{a_0},s_x)$, the clockwise
    offset of their start from $a_0$.
    \item Go through them in that order. Mark $x$ \emph{reachable} if
    $F(a_0,x)$, or if some earlier reachable $y$ has $F(y,x)$, and remember $y$.
    \item As soon as a reachable $x$ also has $F(x,a_0)$, the chain wraps the
    whole circle. Follow the remembered links back to $a_0$ and output the chain.
  \end{itemize}
  \item If no starting lecture works, output ``no valid subset exists''.
\end{enumerate}

\section{Why it is correct}

\emph{Any chain found is valid.} Let the chain be $a_0=x_0,\dots,x_m$ with
$F(x_i,x_{i+1})$ and $F(x_m,a_0)$ and increasing offsets. (Two different lectures
with the same start overlap, so $F$ never links them, and the offsets strictly
increase.) Put each $x_i$ on a line starting at $a_0$: it occupies
$[p_i,p_i+\ell_i)$ where $p_i$ is its offset. Because $x_{i+1}$ does not start
inside $x_i$, we get $p_i+\ell_i\le p_{i+1}$. Likewise $x_m$ must not run past the
start of $a_0$, or it would overlap $a_0$, which condition 1 of $F(x_m,a_0)$
forbids. So the arcs are disjoint, in the order $x_0,\dots,x_m$, and their
neighbours are exactly the pairs we checked. By the claim, the set is valid.

\emph{No valid set is missed.} Suppose a valid $S$ exists. If $|S|=1$, step 2
finds a lecture that works. Otherwise some member $a_0$ of $S$ overlaps
lecture~1, so it is tried in step 4. Listing $S$ clockwise from $a_0$ gives a
chain with increasing offsets, and $F$ holds on all neighbouring pairs by the
claim. The reachability search explores all such chains, so it finds one (maybe a
different valid one), for this $a_0$ or an earlier one.

\section{Complexity}

\begin{itemize}
  \item Overlap table: $O(n^2)$.
  \item Table $F$: $n^2$ pairs, $O(n)$ each, so $O(n^3)$.
  \item One chain search: sorting is $O(n\log n)$, and the reachability loop is
  $O(n^2)$ lookups. At most $n$ starting lectures gives $O(n^3)$.
\end{itemize}
Total: $O(n^3)$ time and $O(n^2)$ space for the two tables. The first brute-force
version tried every subset, which is $O(2^n)$. I did not try to prove
$O(n^3)$ is optimal; it is just the easiest bound to justify.

\section{Testing}

The repository has 11 hand-checked test cases in \texttt{tests/all\_testcases.txt}.
They cover a single lecture, one lecture covering all, identical lectures,
disjoint lectures, separate groups, a chain with a unique answer, touching
lectures, midnight wrap-around, and rings of 4, 5 and 6 lectures. A ring where
each lecture overlaps only its two neighbours has a valid subset exactly when
its size is divisible by 3, which is why the rings of 4 and 5 have no answer and
the ring of 6 does.

\end{document}
